\documentclass[letterpaper,11pt]{article}

\usepackage{latexsym}
\usepackage[empty]{fullpage}
\usepackage{titlesec}
\usepackage[usenames,dvipsnames]{xcolor}
\usepackage{enumitem}
\usepackage[hidelinks]{hyperref}
\usepackage{fancyhdr}
\usepackage[english]{babel}
\usepackage{tabularx}
\usepackage[T1]{fontenc}

%----------FONT OPTIONS----------
\usepackage{charter} % Professional Serif Font

%----------PAGE SETUP----------
\pagestyle{fancy}
\fancyhf{} % clear all header and footer fields
\fancyfoot{}
\renewcommand{\headrulewidth}{0pt}
\renewcommand{\footrulewidth}{0pt}

% Adjust margins
\addtolength{\oddsidemargin}{-0.5in}
\addtolength{\evensidemargin}{-0.5in}
\addtolength{\textwidth}{1in}
\addtolength{\topmargin}{-.5in}
\addtolength{\textheight}{1.0in}

\urlstyle{same}

\raggedbottom
\raggedright
\setlength{\tabcolsep}{0in}

%----------SECTION FORMATTING----------
\titleformat{\section}{
  \vspace{-4pt}\scshape\raggedright\large
}{}{0em}{}[\color{black}\titlerule \vspace{-5pt}]

%----------CUSTOM COMMANDS----------
\newcommand{\resumeItem}[1]{
  \item\small{
    {#1 \vspace{-2pt}}
  }
}

\newcommand{\resumeSubheading}[4]{
  \vspace{-1pt}\item
    \begin{tabular*}{0.97\textwidth}[t]{l@{\extracolsep{\fill}}r}
      \textbf{#1} & #2 \\
      \textit{\small#3} & \textit{\small #4} \\
    \end{tabular*}\vspace{-7pt}
}

\newcommand{\resumeProjectHeading}[2]{
    \item
    \begin{tabular*}{0.97\textwidth}{l@{\extracolsep{\fill}}r}
      \small#1 & #2 \\
    \end{tabular*}\vspace{-7pt}
}

\newcommand{\resumeSubHeadingListStart}{\begin{itemize}[leftmargin=0.15in, label={}]}
\newcommand{\resumeSubHeadingListEnd}{\end{itemize}}
\newcommand{\resumeItemListStart}{\begin{itemize}}
\newcommand{\resumeItemListEnd}{\end{itemize}\vspace{-5pt}}

%----------DYNAMIC STATS SETUP----------
% This connects your resume to the Python script.
% If the file exists, it uses the real numbers. If not, it defaults to placeholders.
\IfFileExists{dynamic_stats.tex}{
    
% AUTOMATICALLY GENERATED FILE
\newcommand{\leetcodeRating}{1893}
\newcommand{\leetcodeSolved}{411}
\newcommand{\codeforcesRating}{1406}
\newcommand{\lastUpdated}{September 20, 2026}

}{
    \newcommand{\leetcodeRating}{0}
    \newcommand{\leetcodeSolved}{0}
    \newcommand{\codeforcesRating}{0}
    \newcommand{\lastUpdated}{Not Updated}
}

%-------------------------------------------
%%%%%%  RESUME STARTS HERE  %%%%%%%%%%%%%%%%%%%%%%%%%%%%

\begin{document}

%----------HEADER----------
\begin{center}
    \textbf{\Huge \scshape Ganji Mahesh} \\ \vspace{5pt}
    \small +91 7032436352 $|$ \href{mailto:maheshganji057@gmail.com}{\underline{maheshganji057@gmail.com}} $|$ 
    \href{https://linkedin.com/in/g-mahesh-781237283}{\underline{g-mahesh-781237283}} $|$
    \href{https://github.com/Mahesh0002}{\underline{github.com/Mahesh0002}}
\end{center}


\section{Education}
  \resumeSubHeadingListStart
    \resumeSubheading
      {Vardhaman College of Engineering}{Hyderabad, Telangana}
      {B.Tech in Computer Science}{2023 -- Present}
      
    \resumeSubheading
      {Sri Chaitanya Junior Kalasala}{Hyderabad, Telangana}
      {Intermediate (MPC)}{2021 -- 2023}
  \resumeSubHeadingListEnd


%-----------EXPERIENCE-----------
\section{Experience}
  \resumeSubHeadingListStart

    \resumeSubheading
      {Rejolt}{Hyderabad, Telangana}
      {Fellowship}{Jul 2025 -- Present}
      \resumeItemListStart
        \resumeItem{Gained hands-on experience with design and prototyping tools like {Figma} to create user-centric interfaces.}
        \resumeItem{Collaborated with the development team to understand the transition from design wireframes to functional code.}
        \resumeItem{Assisted in designing UI components and streamlined user workflows for web applications.}
      \resumeItemListEnd

  \resumeSubHeadingListEnd

%-----------PROJECTS-----------
\section{Projects}
    \resumeSubHeadingListStart

      % PROJECT 1: Car Damage
      \resumeProjectHeading
          {\textbf{Car Damage Classification} $|$ \emph{Python, YOLO-v11, EfficientNet, ViT}}{Aug 2025 -- Jan 2026}
          \resumeItemListStart
            \resumeItem{Engineered a damage assessment pipeline utilizing {YOLO-v11} for precise object localization on vehicle images.}
            \resumeItem{Conducted a comparative analysis between {EfficientNet} and {Vision Transformers} to determine the optimal architecture for damage severity classification.}
            \resumeItem{Benchmarked model performance across categories (minor, moderate, severe), achieving higher accuracy by selecting the best-performing architecture for the final solution.}
          \resumeItemListEnd

      % PROJECT 2: Jōtai
      \resumeProjectHeading
          {\textbf{Jōtai: Advanced Productivity Dashboard} $|$ \emph{React.js, Tailwind CSS}}{Jan 2025 -- May 2025}
          \resumeItemListStart
            \resumeItem{Architected a dynamic dashboard as a technical benchmark to rigorously test and demonstrate proficiency in {React} patterns and state management.}
            \resumeItem{Leveraged {Tailwind CSS} to master utility-first styling, creating a complex, responsive grid layout to evaluate UI rendering performance.}
            \resumeItem{Focused on component modularity and clean code principles to validate personal growth in frontend architecture.}
          \resumeItemListEnd

      % PROJECT 3: Student Community Forum
      \resumeProjectHeading
          {\textbf{Student Collaboration Hub} $|$ \emph{JavaScript, HTML5, CSS3}}{Jul 2024 -- Nov 2024}
          \resumeItemListStart
            \resumeItem{Designed a web-based prototype to unify asynchronous discussions and synchronous video interaction (Zoom-like interface) into a single platform.}
            \resumeItem{Implemented core functionality including user authentication, threaded discussion boards, and search features using vanilla {JavaScript}.}
            \resumeItem{Created a seamless UI to demonstrate the concept of bridging the gap between study forums and live study groups.}
          \resumeItemListEnd

    \resumeSubHeadingListEnd


%-----------TECHNICAL SKILLS-----------
\section{Technical Skills}
 \begin{itemize}[leftmargin=0.15in, label={}]
    \small\item{
     \textbf{Languages}{: Java, Python, C/C++, SQL, JavaScript, HTML/CSS} \\
     \textbf{Machine Learning}{: YOLO, EfficientNet, Transformers, PyTorch, Scikit-learn, Pandas} \\
     \textbf{Web Development}{: React.js, Tailwind CSS, JSP, JDBC, Servlets, Node.js} \\
     \textbf{Tools \& Platforms}{: Git, VS Code, Android Studio, Blender, {Figma}}
    }
 \end{itemize}

 
%-----------ACHIEVEMENTS-----------
% This section now pulls data from dynamic_stats.tex
\section{Achievements}
 \begin{itemize}[leftmargin=0.15in, label=$\bullet$]
    \small
    \item \textbf{LeetCode:} Max Rating \textbf{\leetcodeRating} $|$ Solved \textbf{\leetcodeSolved}+ Problems (Combined)
    \item \textbf{Codeforces:} Max Rating \textbf{\codeforcesRating}
    \item \textbf{English Proficiency:} C1 Level (Lingua Business)
    
    % Optional: Uncomment this line if you want the date visible on the PDF
    % \item \textit{\footnotesize (Stats last auto-updated on \lastUpdated)}
 \end{itemize}

\end{document}